\section{Experiments and Results}

\subsection{Benchmarks and Setup}

We evaluate our memory agent on two long-horizon agent benchmarks:
Terminal-Bench 2.0 and $\tau^2$-Bench. Terminal-Bench 2.0 evaluates autonomous agents
operating in realistic command-line environments, where agents must inspect
files, run commands, edit code, debug failures, and satisfy hidden verifier tests
\citep{merrill2026terminalbench}. In this setting, execution state decays through
local debugging loops: agents may forget earlier requirements, repeat failed
commands, overlook environment observations, or lose track of diagnosed errors.

$\tau^2$-Bench evaluates conversational tool-use agents in dynamic task
environments derived from real-world service domains \citep{yao2024taubench,
barres2025tau2bench}. Unlike purely autonomous terminal tasks, these interactions
require the agent to maintain policy-relevant state across a multi-turn
conversation, coordinate with a user simulator, and use tools to update or verify
the environment. This creates a different source of behavioral state decay:
important facts may be stated by the user, observed through tools, or implied by
domain policies, but later stop influencing the agent's decisions.

Together, these benchmarks allow us to test whether memory  helps
across both autonomous execution and interactive coordination. Terminal-Bench
emphasizes environment grounding, debugging continuity, and procedural memory;
$\tau^2$-Bench emphasizes policy adherence, user-state tracking, and
conversation-level execution state.

\paragraph{Terminal-Bench 2.0.}
We run the Terminal-Bench 2.0 full set using the official Terminus-2-style
terminal agent harness. The benchmark contains 89 tasks; we report results on
the 85 paired tasks for which both baseline and memory runs produced valid
evaluations, excluding four docker failures unrelated to agent behavior.
For each task, the action agent interacts with a containerized terminal
environment and receives a binary pass/fail score from the task verifier. We
report pass@1 averaged across tasks. For Terminal-Bench, pass@1 is the fraction of tasks whose final verifier
passes.

\paragraph{$\tau^2$-Bench.}
We evaluate three domains from the $\tau^2$-Bench framework: airline, retail, and
telecom. We use the base task split, containing 50 airline tasks, 114 retail
tasks, and 114 telecom tasks, for a total of 278 tasks per configuration. Each
episode is a single sampled conversation between the action agent and a user
simulator, with task success determined by the benchmark evaluator. We report
pass@1 for each domain and a task-weighted average across domains.
For $\tau^2$-Bench, a conversation is considered a pass if it satisfies the task evaluator.

\begin{table*}[t]
\centering
\small
\setlength{\tabcolsep}{6pt}
\caption{
Main results on Terminal-Bench and $\tau^2$-Bench. All scores are pass@1
percentages; $\Delta$ is reported in percentage points (pp). The memory agent
is Claude Opus 4.6.
}
\label{tab:main_results}
\begin{tabular}{lllrrrr}
\toprule
Benchmark & Domain / Split & Action Model & $n$
& Baseline & + Memory & $\Delta$ \\
\midrule
Terminal-Bench 2.0 & full set & Sonnet 4.5 & 85
& 37.6\% & \textbf{45.9\%} & \textbf{+8.3 pp} \\
Terminal-Bench 2.0 & full set & Opus 4.6 & 85
& 43.5\% & \textbf{45.9\%} & \textbf{+2.4 pp} \\
\midrule
$\tau^2$-Bench & airline & Sonnet 4.5 & 50
& 68.0\% & \textbf{78.0\%} & \textbf{+10.0 pp} \\
$\tau^2$-Bench & retail & Sonnet 4.5 & 114
& 49.1\% & \textbf{58.8\%} & \textbf{+9.6 pp} \\
$\tau^2$-Bench & telecom & Sonnet 4.5 & 114
& 55.3\% & \textbf{57.9\%} & \textbf{+2.6 pp} \\
\cmidrule(lr){2-7}
$\tau^2$-Bench & task-weighted avg. & Sonnet 4.5 & 278
& 55.0\% & \textbf{61.8\%} & \textbf{+6.8 pp} \\
\midrule
$\tau^2$-Bench & airline & Opus 4.6 & 50
& 76.0\% & 76.0\% & +0.0 pp \\
$\tau^2$-Bench & retail & Opus 4.6 & 114
& 64.9\% & \textbf{69.3\%} & \textbf{+4.4 pp} \\
$\tau^2$-Bench & telecom & Opus 4.6 & 114
& 63.2\% & \textbf{64.9\%} & \textbf{+1.8 pp} \\
\cmidrule(lr){2-7}
$\tau^2$-Bench & task-weighted avg. & Opus 4.6 & 278
& 66.2\% & \textbf{68.7\%} & \textbf{+2.5 pp} \\
\bottomrule
\end{tabular}
\end{table*}
\paragraph{Agents and memory configuration.}
We evaluate two action-agent strengths: Claude Sonnet 4.5 and Claude Opus 4.6.
Unless otherwise stated, the memory agent is Claude Opus 4.6. The action agent is
kept unchanged across baseline and memory-enabled runs. In the memory condition,
the memory agent runs at the first step and then at every subsequent step. It observes the task description, a recent trajectory
window of $k$ ($k$=8) messages, and the current memory bank. It first updates the
memory bank, then emits either a context injection or
no intervention. If a reminder is emitted, it is injected into the next action-agent call as transient memory context.


\subsection{Main Results}
Table~\ref{tab:main_results} reports the main results. Memory intervention improves pass@1 on both benchmarks and across both action-agent strengths: on Terminal-Bench 2.0, Sonnet 4.5 gains +8.3 pp ($37.6\% \to 45.9\%$) and Opus 4.6 gains +2.4 pp ($43.5\% \to 45.9\%$); on $\tau^2$-Bench, Sonnet 4.5 gains +6.8 pp on the task-weighted average ($55.0\% \to 61.8\%$) and Opus 4.6 gains +2.5 pp ($66.2\% \to 68.7\%$). Gains are larger for the weaker action agent but do not disappear for the stronger one, indicating that the benefit is not merely compensating for limited capacity. Across $\tau^2$-Bench domains, airline and retail show the largest Sonnet lifts (+10.0 and +9.7 pp) while telecom is smaller (+2.6 pp), consistent with our framing of memory as a domain-sensitive intervention policy rather than a fixed summarization mechanism.

\begin{table*}[t]
\centering
\small
\setlength{\tabcolsep}{4pt}
\caption{
Ablations on $\tau^2$-Bench with Sonnet 4.5 as the action agent, Opus 4.6 as the memory agent for all experiments. Scores are
pass@1 percentages. Macro averages domains equally, micro is task-weighted.
}
\label{tab:ablation_all}
\resizebox{\textwidth}{!}{%
\begin{tabular}{lllccccc}
\toprule
 & & & \multicolumn{3}{c}{$\tau^2$-Bench pass@1 (\%)} & \multicolumn{2}{c}{Average} \\
\cmidrule(lr){4-6} \cmidrule(lr){7-8}
Variant & Phase 1 (memory) & Phase 2 (intervention) & Airline & Retail & Telecom & Macro & Micro \\
\midrule
Sonnet 4.5 baseline
& --- & ---
& 68.0 & 49.1 & 55.3 & 57.5 & 55.0 \\
\midrule
\textbf{+ Full memory agent (ours)}
& bank management
& selective reminder / silence
& \textbf{\imp{78.0}} & \imp{57.0} & \imp{57.9}
& \textbf{\imp{64.3}} & \imp{61.2} \\
+ Full-bank context
& bank management
& expose full bank
& \imp{74.0} & \imp{52.6} & \imp{57.9}
& \imp{61.5} & \imp{58.6} \\
+ Always inject
& bank management
& forced reminder every step
& \imp{72.0} & \imp{58.8} & \imp{59.6}
& \imp{63.5} & \textbf{\imp{61.5}} \\
+ Injection-only (no bank)
& skipped
& selective guidance / silence
& \reg{62.0} & \imp{54.4} & \textbf{\imp{66.7}}
& \imp{61.0} & \imp{60.8} \\
+ Mem0 \citep{mem0ai2026mem0}
& Mem0 ADD
& vector+BM25 top-10
& 68.0 & \textbf{\imp{59.6}} & \imp{58.8}
& \imp{62.1} & \imp{60.8} \\
\bottomrule
\end{tabular}%
}
\end{table*}


\subsection{Ablations and Comparisons}

We ablate the memory agent on $\tau^2$-Bench to test which parts of the architecture are responsible for the gains. Our memory design has two agentic components. Phase~1 performs memory management: it maintains a persistent memory bank of user facts, tool observations, procedural evidence, and open subgoals. Phase~2 performs intervention: it decides whether any remembered state should be reactivated in the next action-agent call, and if so, generates a concise memory-grounded reminder. The ablations remove one capability at a time, allowing us to compare our design against passive context augmentation, always-on memory injection, advisor-style guidance.

\paragraph{Ablation variants.}
\textit{Full-bank context} keeps Phase~1 but removes Phase~2 intervention
selection. The memory agent still performs agentic memory management, but the
entire memory bank is exposed to the action agent at every step. This resembles
passive context augmentation: the system maintains memory, but does not decide
which memory item is currently relevant or whether it should intervene in the
control loop. This tests whether the benefit comes merely from making historical
state visible.

\textit{Always inject} keeps both the memory bank and the generative Phase~2
reminder, but removes the silence action. Unlike embedding-based
retrieval-augmented prompting, which retrieves fixed records from a store
\citep{lewis2020rag}, this variant still uses the memory model to synthesize a
reminder from the maintained bank. However, it must inject at every step. This
tests whether the no-op decision is only an efficiency optimization or an
essential part of intervention calibration.

\textit{Injection-only (no bank)} removes persistent memory while retaining the
Phase~2 intervention model. This is analogous to advisor-style systems, where an
auxiliary model observes the trajectory and provides natural-language guidance
to an executor \citep{asawa2025advisor,anthropic2026advisor}. Unlike the full
memory agent, this variant has no maintained memory bank, so its guidance cannot
be grounded in persistent execution state.

\textit{Mem0} compares against a general persistent memory layer
\citep{mem0ai2026mem0}. In this variant, memories are written through Mem0's
ADD interface and retrieved with its vector/BM25 search, with the top retrieved
items returned to the action agent. This isolates the difference between
production memory retrieval and our formulation of memory as a selective
intervention policy.


Table~\ref{tab:ablation_all} shows that memory-style variants generally improve
over the Sonnet baseline, but the improvements are not equally robust across
domains. The full two-phase memory agent achieves the highest macro-average,
improving all three domains and giving the largest gain on airline. Since macro
averages domains equally, this indicates that our architecture gives the most
balanced improvement across heterogeneous task types.

The ablations show that each removed capability weakens robustness in a
different way. Full-bank context improves over the baseline but trails the full
system by 2.8 macro points and 2.6 micro points, suggesting that maintaining a
memory bank is useful but exposing all remembered state is not enough; the memory
agent must select what is currently behaviorally relevant. Always inject is
competitive and slightly leads on micro-average by 0.3 points, but this gap is
within expected run variance and disappears on the domain-balanced macro-average,
where selective silence is better, especially on airline. Injection-only guidance helps in some domains, especially
telecom, but is less stable without persistent memory and even hurts airline
relative to the baseline. Together, these results suggest that neither passive
memory exposure, always-on reminders, nor generic auxiliary guidance is
sufficient: the most balanced gains come from combining maintained execution-state
memory with a selective intervention policy.

Mem0 improves the average score, confirming that persistent
memory retrieval is useful. However, it does not improve airline over the
baseline and falls short of our full system on macro-average. This highlights
the distinction between retrieving memories and deciding whether retrieved state
should intervene. A general memory layer can return relevant records, but it
does not explicitly model when a memory should enter the loop or how it
should be phrased as a targeted reminder.

Overall, the ablations support the two-phase design. Agentic memory management
alone is helpful but insufficient; advisor-like intervention without memory is
not reliably grounded; and always-on injection can be competitive but less
balanced. The full memory agent performs best on macro-average because it
combines maintained execution state with a selective intervention policy that can
choose both what to say and when to remain silent.


\subsection{Qualitative Analysis}

We qualitatively inspect memory-enabled trajectories to understand when memory
injection changes agent behavior. Across both benchmarks, the dominant pattern
is not generic summarization. Memory helps when it reactivates a specific piece
of execution state at the moment the action agent is about to ignore it. The
type of state differs by benchmark: Terminal-Bench emphasizes environment
observations, debugging continuity, and failure diagnoses, while $\tau^2$-Bench
emphasizes policy clauses, user/tool state, and procedural sequencing.
Table~\ref{tab:qual_mechanisms} summarizes the recurring mechanisms we observe;
the rest of this section discusses each benchmark in turn.

\begin{table*}[t]
\centering
\small
\setlength{\tabcolsep}{5pt}
\caption{
Representative mechanisms observed in qualitative analysis. Memory interventions
help when they reactivate execution state that the action agent previously saw
but no longer uses in its next decision.
}
\label{tab:qual_mechanisms}
\begin{tabularx}{\textwidth}{l X X}
\toprule
Mechanism & Typical memory content & Representative examples \\
\midrule
Requirement / policy reactivation
& Task or domain rule about an allowed action
& $\tau^2$ airline compensation, retail modification rules \\
Environment grounding
& Runtime facts, paths, tool limitations, system quirks
& Terminal-Bench Git server setup, ARS file-write failure \\
Failure-loop avoidance
& Previous attempts and why they failed
& Adaptive rejection sampling, telecom diagnostic retries \\
Diagnostic carryover
& Root cause of a bug or negative signal
& Regex edge cases, SQLite gcov configuration \\
Progress / entity tracking
& Which user, order, line, branch, or subgoal is active
& $\tau^2$ telecom line lookup, retail authentication state \\
\bottomrule
\end{tabularx}
\end{table*}

\paragraph{Terminal-Bench: maintaining debugging continuity.}
In Terminal-Bench, memory is most useful during iterative debugging, where the
action agent accumulates observations but later loses track of which ones should
constrain the next command. The relevant memory is not a single type of
fact. In \texttt{regex-log}, memory reactivates both requirements and diagnoses:
it points out that the current regex violates the task's boundary condition,
misses single-digit IPv4 octets. In \texttt{adaptive-rejection-sampler}, memory
tracks repeated failed file edits and later surfaces an environment-specific
workaround. These cases suggest that memory improves Terminal-Bench performance
by preserving continuity across debugging iterations: what was tried, what was
observed, what failed, and what still constrains the solution.

\paragraph{$\tau^2$-Bench: reactivating policy and interaction state.}
In $\tau^2$-Bench, the useful execution state is often policy- or
interaction-specific. Airline and retail tasks frequently require the agent to
apply a rule many turns after the relevant fact was observed. For example, in
one airline case, the user claimed Gold status, but tool output showed that the
user was a Regular member. The baseline issued compensation based on the user's
claim, while the memory-enabled agent reminded the action agent to rely on
verified records. In another case, memory prevented an invalid modification of a
basic-economy flight by reactivating the policy clause that such flights cannot
be modified. Successful interventions often occur immediately before a
state-changing tool call, reminding the agent of authentication requirements,
eligibility conditions, one-shot tool limits, or policy clauses not enforced.

Overall, this qualitative analysis aligns with the main results: memory
interventions help more often than they hurt. Successful interventions are
specific, grounded, and timely; they restore the causal influence of execution
state over the next action. The remaining failures are primarily calibration
errors rather than failures of memory storage. Occasionally, the memory agent
surfaces a speculative inference with too much confidence, repeats information
that the action agent already knows, or raises a plausible but unnecessary
concern that causes extra verification. These cases motivate training the memory
agent to better decide when to remain silent.

\subsection{Training Open-Weight Memory Agents}
\label{sec:training}

We conduct an early exploration of whether the memory-intervention policy can be
trained on an open-weight model. The prompted Opus memory agent improves
long-horizon agents, but it adds a frontier-model call at each memory step and
its intervention calibration is imperfect. We therefore train Qwen3.5-27B as a
memory agent while keeping the Qwen3.5-122B-A10B action agent frozen.

We train on SETA, a collection of executable terminal-agent tasks with verifier
rewards \citep{seta}. SETA is used for supervised fine-tuning and reinforcement
learning, and Terminal-Bench 2.0 is held out to test transfer. SFT distills
prompted-memory trajectories, supervising both Phase~1 memory-bank operations
and Phase~2 intervention decisions. GRPO then further optimizes intervention
calibration. Because task-level verifier rewards are sparse over many memory
calls, we focus updates on pivot turns that labeled offline rollouts identify as likely to affect downstream success.

\begin{table}[t]
\centering
\small
\caption{
Training and transfer of the open-weight memory agent. The action agent is a
frozen Qwen3.5-122B-A10B; the trainable memory agent is Qwen3.5-27B. SETA
validation reports average verifier reward; Terminal-Bench reports pass@1 on
the 85-task.
}
\label{tab:trained_memory}

\begin{subtable}[t]{0.49\textwidth}
\centering
\setlength{\tabcolsep}{5pt}
\caption{SETA validation.}
\label{tab:seta_trained_memory}
\begin{adjustbox}{max width=\linewidth}
\begin{tabular}{lccc}
\toprule
Setup & Avg. reward & Solved & $\Delta$ \\
\midrule
Action only, no memory
& 0.709 & 56 & --- \\
+ Qwen3.5-27B base memory
& 0.693 & 54 & -0.016 \\
+ SFT memory
& 0.720 & 58 & +0.011 \\
+ GRPO memory
& \textbf{0.734} & \textbf{58} & \textbf{+0.025} \\
\bottomrule
\end{tabular}
\end{adjustbox}
\end{subtable}
\hfill
\begin{subtable}[t]{0.49\textwidth}
\centering
\setlength{\tabcolsep}{6pt}
\caption{Transfer to Terminal-Bench 2.0.}
\label{tab:trained_memory_terminalbench}
\begin{adjustbox}{max width=\linewidth}
\begin{tabular}{lccc}
\toprule
Setup & $n$ & Pass@1 & $\Delta$ \\
\midrule
Qwen3.5-122B-A10B action only
& 85 & 37.6\% & --- \\
+ trained Qwen3.5-27B memory
& 85 & \textbf{41.1\%} & \textbf{+3.5 pp} \\
\bottomrule
\end{tabular}
\end{adjustbox}
\end{subtable}
\end{table}

Table~\ref{tab:seta_trained_memory} shows that the memory policy is learnable
but requires calibration. An untrained 27B memory agent hurts performance,
reducing SETA reward from 0.709 to 0.693. SFT recovers this loss and improves
reward to 0.720, while GRPO further raises it to 0.734. On held-out
Terminal-Bench 2.0, the SETA-trained memory agent improves the frozen action
agent from 37.6\% to 41.1\% pass@1. These results suggest that memory
intervention can be distilled into an open-weight model and that RL improves the
decision of when remembered state should enter the control loop.